% =====================================================================================
%  Complexity and measured runtime. Numbers from code/results_evaluation/runtime.csv
%  and complexity.csv. Timing from an instrumented run: 160 solves, 20 per variant.
% =====================================================================================

\subsection{Complexity and measured cost}
\label{sec:complexity}

\paragraph{Asymptotic cost.}
Write $n$ for the number of work items in an arrival batch, $m$ for the number of eligible
groups, $N$ for population size, $M$ for the number of objectives and $H$ for the number of
reference directions. The operational baselines are single-pass: random, round robin and FIFO
are $O(n)$, least-loaded is $O(n \log m)$ with a heap over group loads, and greedy is $O(nm)$
because it scans every eligible group for every item. The evolutionary methods are dominated
per generation by non-dominated sorting at $O(M N^2)$ \citep{deb2002fast}; NSGA-III adds
association to reference directions, giving $O(M N^2 + N H)$ \citep{debjain2014evolutionary}.
The repair operator is $O(nm)$ in the worst case and is bounded at six passes, so it does not
change the order.

The proposed modifications do not alter the asymptotic cost. Priority-aware initialisation is
a one-off $O(nm)$ greedy construction over a fraction of the initial population. Adaptive
crossover and the mutation schedule each add an $O(N n)$ diversity measurement per generation,
dominated by the $O(M N^2)$ sorting term.

\paragraph{Measured cost.}
Asymptotic equality does not imply equal wall time, so Table~\ref{tab:runtime} reports
measured per-solve time from an instrumented run of 160 solves, 20 per variant, at 150
generations and population 60 on CPU with no GPU. Timing is captured at the point the search
executes rather than derived from log timestamps.

\begin{table}[htbp]
\centering
\caption{Measured wall time per solve. 20 solves per variant, 150 generations, population 60,
single CPU. Percentages are relative to standard NSGA-II.}
\label{tab:runtime}
%% GENERATED by code/figures/make_tables.py. Do not edit by hand.
%% The caption, label and float live in the manuscript; only the numbers are here.
\begin{tabular}{lrrrr}
\toprule
Policy & Mean s/solve & SD & Solves & Percent vs NSGA-II \\
\midrule
Standard NSGA-II               & 0.89 & 0.02 & 20 & 0.0 \\
NSGA-III                       & 0.99 & 0.01 & 20 & 12.2 \\
PAA-NSGA-II                    & 0.96 & 0.01 & 20 & 8.1 \\
PAA without priority init      & 0.96 & 0.01 & 20 & 7.8 \\
PAA without adaptive crossover & 0.91 & 0.02 & 20 & 2.4 \\
PAA without SLA mutation       & 0.98 & 0.02 & 20 & 10.1 \\
PAA with adaptive mutation     & 0.92 & 0.01 & 20 & 3.5 \\
NSGA-III + priority init       & 1.01 & 0.02 & 20 & 14.3 \\
\bottomrule
\end{tabular}


\end{table}

\begin{figure}[htbp]
\centering
\includegraphics[width=0.72\textwidth]{fig09_runtime.pdf}
\caption{Measured wall time per solve, 20 solves per variant, single CPU. Every evolutionary
variant lies within \SI{14}{\percent} of standard NSGA-II, so on this problem the choice
between them is not a cost decision.}
\label{fig:runtime}
\end{figure}

Three observations follow.

The spread is narrow. Every evolutionary variant lies within \SI{14}{\percent} of standard
NSGA-II, and at \num{0.89} to \num{1.01} seconds per batch of 40 items, with a standard
deviation under \num{0.02} across twenty solves each, all are far inside the latency budget of an
assignment decision a human would otherwise take in seconds. These timings were measured on
an otherwise idle machine. An earlier set, taken while four other jobs were running, read two
to three times higher and with twenty times the variance, and is recorded here because a
runtime measured under contention is not a runtime. These timings were measured on a workstation that was not
otherwise idle, so they are reported as indicative of the ordering and not as benchmarks.

Priority-aware initialisation buys search quality and is paid for in time. Adding it to
NSGA-III raises hypervolume from \num{0.0837} to \num{0.0862}, and raises mean solve time
from \num{1.80} to \num{3.19} seconds, which is \SI{77}{\percent} longer.

An earlier draft reported this as a component that pays for itself twice, raising quality
and cutting time together. That was measured before the hypervolume reference point was
fixed and before the timing instrumentation was corrected, and it does not survive either.
Constructing a priority-aware starting population costs real work and the run does not win
it back. The trade is still favourable at this scale, because the extra \num{1.38} seconds sits
well inside the budget of a decision a human would take in seconds, but it is a trade and
it is reported as one.

The two modifications that do not help do not cost much either. Removing adaptive crossover or
SLA-aware mutation saves roughly \SI{5}{\percent} of solve time while changing hypervolume by
less than \num{0.001}. They are close to free and close to useless, which is the honest
summary and the reason the ablation is reported per component.
